\chapter{Conclusiones y trabajo futuro}

\section{Conclusiones}

\subsection{Síntesis del problema y el enfoque adoptado}

Este trabajo aborda la predicción dinámica del resultado de partidos de fútbol como un problema de clasificación multiclase supervisada. A diferencia de los sistemas de predicción prepartido que dominan la literatura, el sistema desarrollado estima la probabilidad de victoria local, empate y victoria visitante a partir de métricas acumuladas durante el propio partido, actualizando la predicción en cada minuto del tiempo reglamentario.

El fútbol es un deporte de baja puntuación y alta variabilidad entre encuentros, en el que un único gol puede alterar radicalmente el desenlace. Estas características hacen que la predicción de resultados sea reconocidamente difícil, incluso para enfoques basados en aprendizaje automático con datos extensos. El trabajo se posiciona en un punto intermedio entre los modelos puramente prepartido, que no incorporan la evolución del juego, y las arquitecturas secuenciales profundas, que modelan los eventos individuales con mayor complejidad pero menor interpretabilidad. En su lugar, se opta por transformar el registro de eventos de StatsBomb Open Data en un conjunto de métricas temporales por minuto, construyendo un dataset de modelado con 48 variables sobre 311.760 observaciones que sirve de entrada a clasificadores de ensamblado.

\subsection{Validez de la solución}

Los resultados obtenidos confirman que el enfoque es técnicamente válido. Los tres modelos entrenados superan con claridad la línea base de predicción sistemática de la clase mayoritaria, que alcanza un 42,9\,\% de precisión, con mejoras de entre 21 y 22,5 puntos porcentuales. El mejor modelo, XGBoost, obtiene un 65,4\,\% de precisión global y una pérdida logarítmica de 0,753, situándose en el mismo rango que los modelos prepartido del estado del arte, a pesar de operar con información parcial del partido en curso.

El análisis temporal constituye el argumento más sólido en favor de la validez del sistema: la precisión crece de forma consistente a lo largo del partido, superando el umbral de referencia del 70\,\% a partir del minuto 60 en el caso de Random Forest. Este comportamiento es coherente con la dinámica real del fútbol, donde la información acumulada en la segunda mitad del partido es significativamente más informativa que la de los primeros minutos. Un sistema que mejora su fiabilidad a medida que avanza el encuentro es precisamente lo que se espera de un predictor dinámico intrapartido.

La validez del enfoque no se limita al rendimiento de los modelos. El prototipo desarrollado ha sido desplegado como aplicación web en Streamlit Community Cloud y permite visualizar la predicción en cualquier minuto de un partido del conjunto de test, observar la evolución de las probabilidades a lo largo del encuentro y explorar escenarios hipotéticos mediante el simulador de estado de partido. La existencia de este prototipo funcional demuestra que los modelos entrenados son integrables en un sistema de uso real y que el enfoque tiene proyección más allá del ámbito experimental.

\subsection{Contribuciones del trabajo y relación con los objetivos}

El trabajo ha producido cuatro contribuciones concretas que dan respuesta a los objetivos planteados en el capítulo de Objetivos y metodología.

La primera es la construcción de un dataset de modelado específico para la predicción dinámica intrapartido. A partir de 12,2 millones de eventos de StatsBomb Open Data correspondientes a 3.464 partidos de 21 competiciones, se ha generado un dataset tabular de 311.760 filas y 48 variables que captura el estado del partido en cada minuto del tiempo reglamentario. Este dataset incluye estadísticas acumuladas, variables de ventana deslizante de 15 minutos y variables derivadas que codifican diferenciales de marcador, posesión, producción ofensiva y calidad de los tiros. Su construcción responde directamente a los objetivos de recopilación y tratamiento de datos y de ingeniería de características.

La segunda contribución es la comparativa experimental entre tres clasificadores multiclase, Random Forest, XGBoost y SVM lineal, bajo un protocolo riguroso de validación temporal a nivel de partido que evita la fuga de información. La búsqueda de hiperparámetros mediante \texttt{RandomizedSearchCV} con partición interna por \texttt{match\_id} garantiza que los resultados reflejan la capacidad de generalización real de los modelos a partidos no vistos durante el entrenamiento. Esta comparativa cubre el objetivo de entrenamiento y comparación de clasificadores, y la exclusión justificada del LSTM como modelo alternativo queda respaldada por los resultados obtenidos con los métodos de ensamblado.

La tercera contribución es el marco de evaluación temporal, que va más allá de las métricas globales habituales y analiza el rendimiento de los modelos en seis puntos del partido (minutos 15, 30, 45, 60, 75 y 90). Este análisis permite determinar a partir de qué momento del encuentro el sistema resulta fiable y constituye una metodología evaluativa adecuada para cualquier sistema de predicción dinámica intrapartido. Responde al objetivo de evaluación mediante métricas cuantitativas y análisis temporal.

La cuarta contribución es el prototipo de aplicación interactiva, accesible públicamente en Streamlit Community Cloud. El prototipo integra los modelos XGBoost, Random Forest y SVM lineal en una interfaz que permite explorar partidos del conjunto de test con visualización de la línea de tiempo de eventos y evolución de probabilidades, así como un simulador que acepta el estado del partido como entrada directa. Esto da cumplimiento al objetivo de integración del modelo en una solución interactiva.

El grado de cumplimiento detallado de cada objetivo se recoge en la tabla~\ref{tab:cumplimiento}. La única desviación significativa respecto a la planificación inicial es el incumplimiento del umbral del 70\,\% de precisión global, que no se alcanza con ninguno de los modelos sobre el conjunto de test completo. No obstante, como se argumenta en el capítulo de Resultados, este objetivo se cumple de forma temporal a partir del minuto 60, que es el tramo del partido donde la predicción dinámica resulta más útil en la práctica.

\section{Líneas de trabajo futuro}

Los resultados obtenidos y las limitaciones identificadas abren varias líneas de trabajo que podrían aportar valor añadido al sistema desarrollado.

La extensión más directa es la integración con datos en tiempo real. El sistema actual opera en modo simulación, consumiendo métricas precalculadas de partidos ya disputados. Conectar el pipeline de ingeniería de características a una fuente de eventos en vivo permitiría realizar predicciones durante partidos en curso, que es el escenario de aplicación más valioso para los contextos de \textit{broadcasting} y apuestas \textit{in-play}. Esta extensión requiere replicar el pipeline de construcción de variables en un entorno de baja latencia y establecer una conexión con un proveedor de datos en tiempo real.

En segundo lugar, el simulador de la aplicación actualmente acepta como entrada estadísticas agregadas del partido. Una mejora natural sería implementar un simulador por eventos, en el que el usuario pudiera añadir acciones individuales (gol, tiro, tarjeta…) para un equipo y un minuto determinados, y el sistema recalculara automáticamente las 48 variables y actualizara la predicción. Esta funcionalidad, descartada en este trabajo por la complejidad adicional que implica replicar el pipeline de ingeniería de características dentro de la aplicación, acercaría el simulador a una herramienta de análisis táctico real.

La tercera línea consiste en explorar modelos secuenciales que aprovechen la naturaleza temporal de los datos. Los resultados obtenidos validan que los eventos de StatsBomb contienen señal predictiva suficiente; modelos como LSTM con mecanismos de atención \citep{zhang2022lstm} o arquitecturas Transformer entrenadas sobre secuencias de eventos \citep{simpson2022seq2event} podrían capturar dependencias temporales que la agregación minuto a minuto no representa explícitamente. El marco de evaluación temporal desarrollado en este trabajo podría reutilizarse directamente para comparar estos enfoques con los modelos de ensamblado bajo las mismas condiciones.

Una cuarta línea es la mejora específica en la predicción de empates. Con un \textit{recall} máximo del 52\,\% en el mejor de los modelos, el empate sigue siendo la clase de mayor dificultad. Técnicas como el remuestreo sintético de la clase minoritaria (SMOTE), la reformulación del problema como clasificación ordinal aprovechando el orden natural entre clases, o la definición de umbrales de decisión asimétricos por clase podrían mejorar este comportamiento sin comprometer el rendimiento global.

Finalmente, la ampliación del dataset a otras competiciones reduciría el sesgo hacia el fútbol profesional de élite que introduce la fuente de datos actual. StatsBomb Open Data cubre principalmente ligas europeas de primer nivel, con la Liga española como la más representada (25,1\,\% de los partidos). Incorporar datos de ligas de niveles inferiores, competiciones de otros continentes o un mayor número de temporadas recientes mejoraría la generalización del modelo a contextos de juego más variados y aumentaría la representatividad del conjunto de test.
